
\section{Related Work}
\noindent \methodname builds on recent progress in generative 3D asset creation, but introduces a new form of control. Rather than improving unconditional or globally conditioned generation, our goal is to let users specify where geometric and appearance constraints should apply. 

We therefore review prior work along four directions: 3D asset generation, spatial control for 3D generation, appearance control across 3D representations, and part- or primitive-based editing.

\subsection{3D Asset Generation}
\label{sec:rw_3d_generation}
The field of 3D generation has experienced quick progress in both output quality and representation flexibility. Early diffusion-based 3D generators operated directly in the original 3D space and were limited in efficiency and output type~\cite{nichol2022pointegenerating3dpoint}. More recent approaches moved generation into compact latent spaces, yielding substantial gains in quality and scalability~\cite{NEURIPS2022_40e56dab,jun2023shapegeneratingconditional3d}. Recent asset generators such as CLAY, Hunyuan3D, SAM 3D, and TRELLIS further improve fidelity through large-scale training and richer 3D representations ~\cite{10.1145/3658146,hunyuan3d22025tencent,Chen_2026_CVPR,Xiang_2025_CVPR}. TRELLIS is particularly relevant to our work, as it separates sparse structure generation from structured-latent appearance generation and supports decoding to meshes, radiance fields, or Gaussian splats. While these models provide strong text or image conditioned generation, they do not offer explicit part-wise control over where geometric and appearance conditions should apply.


\subsection{Spatially Controllable 3D Generation} 
\label{sec:rw_spatial_control}

Spatially controllable generation shifts conditioning from ambiguous text or image prompts to explicit 3D inputs, such as voxels, primitives, proxy shapes, or meshes. Existing methods can be grouped into training-based and test-time approaches. 

Training-based methods modify or fine-tune generative models to accept geometric conditions, for example through voxel conditioning, primitive-based layouts, or mesh/proxy inputs~\cite{NEURIPS2022_40e56dab, 10.1145/3641519.3657461}. These approaches can provide direct spatial control, but require additional training and are often tied to the conditioning format and data distribution used during training. 

Test-time methods avoid training the full generator, but introduce other trade-offs. Optimization-based approaches project 3D constraints into rendered views or optimize an explicit 3D representation under geometric guidance~\cite{Metzer_2023_CVPR,10.1145/3641519.3657425}. 

SpaceControl is closest to our work: it directly injects user-specified 3D geometry into the latent space of a pretrained 3D generator and exposes a control strength that trades geometric fidelity for generative realism~\cite{fedele2025spacecontrol}. However, this trade-off is applied globally. A single control strength cannot express common design intent where some regions should remain faithful to the input geometric control while others should be completed or reshaped by the generative prior. 

\subsection{Appearance and Texture Control for 3D Assets} 
\label{sec:rw_appearance_control}

Appearance control for 3D assets has been studied across several 3D representations. Mesh-based methods such as Text2Mesh, TEXTure, EASI-Tex, and TriTex use text prompts, image diffusion priors, or learned texture fields to generate, edit, or transfer appearance on given 3D geometry~\cite{text2mesh, 10.1145/3588432.3591503,perla2024easitex, Cohen-Bar_2025_CVPR}. Other methods extend stylization to implicit or point-based representations, including radiance fields and Gaussian splats~\cite{liu2023stylerf,liu2023stylegaussian,10.1007/978-3-031-85187-2_15}. These approaches can produce high-quality appearance edits, but they are typically tied to a specific output representation or operate on an existing asset rather than controlling the appearance stage of a generative 3D pipeline.

GuideFlow3D is closest to our appearance-control setting, as it performs training-free guidance of a pretrained rectified-flow model and introduces a self-similarity loss for robust appearance transfer~\cite{sayandsarkar_2025_guideflow3d}.  However, its goal is to transfer an overall appearance specification to a given 3D asset. Rather than providing a direct user interface that pairs specific object parts with target cues, current approaches rely on indirect global conditioning (such as text prompts or reference images) to dictate local appearance.


\subsection{Primitive-Based and Part-Aware Editing}
\label{sec:rw_part_primitive_editing}

A complementary line of work aims to make 3D assets editable by exposing an intermediate structure in the form of primitives, parts, graphs, or programs. 

Primitive-abstraction methods approximate shapes with compact geometric elements, ranging from volumetric primitives to superquadrics and learned neural primitives~\cite{abstractionTulsiani17,Paschalidou2019CVPR,Paschalidou2021CVPR}. These decompositions provide interpretable manipulation handles, but are typically designed for abstraction or reconstruction and remain tied to the chosen primitive family or learned parser.


Part-aware generative models make object structure explicit during generation. StructureNet and ShapeAssembly represent shapes through hierarchical part graphs or executable programs over cuboid proxies~\cite{MoGuerreroEtAl:StructureNet:SiggraphAsia:2019,jones2020shapeAssembly}. Later neural approaches build part-level manipulation into implicit shapes, local radiance fields, or diffusion-based point and latent representations, enabling operations such as part transformation, recombination, completion, and mixing~\cite{hertz2022spaghetti,Tertikas2023CVPR,koo2023salad,nakayama2023difffacto}.

Recent methods extend this idea toward more realistic and higher-fidelity asset generation: PartGen reconstructs objects as meaningful editable parts, while CompoSE, the closest prior in this line, uses coarse geometric primitives such as bounding boxes as part controls to synthesize part-separated objects for localized editing~\cite{Chen2025PartGen,slim2026composecompositionalsynthesisediting}.

These methods demonstrate the value of explicit part structure, but typically require the generated asset itself to be represented, generated, or reconstructed as a set of parts.

% Part-aware generative models extend this idea by explicitly modeling object structure. StructureNet represents shapes as hierarchical graphs of semantic parts and relations, enabling structure-aware generation, interpolation, and editing \cite{MoGuerreroEtAl:StructureNet:SiggraphAsia:2019}. ShapeAssembly instead generates procedural programs whose execution produces hierarchical assemblies of cuboid part proxies, combining the editability of programs with learned shape priors
% \cite{jones2020shapeAssembly}. 

% More recent neural generative models build part-level control directly into implicit or point-based representations.
% SPAGHETTI disentangles implicit shapes into controllable shape segments and supports operations such as transforming, interpolating, and recombining parts without explicit part supervision \cite{hertz2022spaghetti}. PartNeRF generates
% objects as collections of locally defined radiance fields, enabling part transformations and part mixing while also modeling appearance \cite{Tertikas2023CVPR}. Diffusion-based approaches such as SALAD and DiffFacto further show that part-level latent representations can support completion, mixing, interpolation, and transformation edits within generative 3D models \cite{koo2023salad,nakayama2023difffacto}.

% Recent work has also begun to revisit part-aware control in the context of modern high-fidelity 3D asset generation. PartGen converts text-, image-, or 3D-conditioned inputs into objects composed of meaningful reconstructed parts,allowing downstream part-level editing of otherwise monolithic generated assets \cite{Chen2025PartGen}. CompoSE is especially close in spirit to spatial part-level control, as it takes coarse geometric primitives such as bounding boxes as part controls and synthesizes part-separated objects that support localized edits \cite{slim2026composecompositionalsynthesisediting}. These methods demonstrate the value of explicit part structure for controllability, but typically require the output
% asset itself to be represented, generated, or reconstructed as a set of parts.

\paragraph{Positioning.}
Taken together, prior work provides high-quality 3D asset generation, spatial conditioning, appearance editing, or explicit part-level editability, but these capabilities remain separate. General asset generators offer strong priors but only global text- or image-level control; spatial-control methods typically impose a global geometry or depend on a fixed conditioning format; texture-editing methods are often post-hoc and tied to a specific 3D representation; and part-aware methods usually require the output asset itself to be generated or reconstructed as an explicit set of parts. \methodname addresses this gap by using user-provided primitives and part assignments as localized control signals for a pretrained 3D asset generator, rather than as a required output representation. This enables users to enforce geometric and appearance constraints only where needed, while allowing unconstrained regions to be completed by the generative prior.
